\begin{center}
\textbf{ABSTRACT}
\end{center}

$\!$\\
\noindent
This work presents the development of an analytical pipeline integrated with a decision-support dashboard for the management of the technological assets of the Hospital das Clinicas of the Federal University of Pernambuco. The study addresses resource allocation under budget constraints, focusing on the replacement of obsolete medical equipment. The proposed methodology includes the automation of extraction, transformation, and loading stages to clean and consolidate operational data, followed by the definition of a failure-criticality prioritization index. Based on this index, the tool implements a budget-distribution utility that, given a financial cap set by the manager, automatically selects the set of equipment to replace, maximizing the aggregate impact of the investment under the budget constraint, thereby translating the multicriteria prioritization into an objective resource-allocation recommendation. Applied to real asset data, this utility turns the prioritized list into an actionable replacement plan, indicating not only \textit{which} equipment to prioritize but \textit{which} to actually replace given the available budget. The contribution is to provide the Clinical Engineering department with an analytical tool to support strategic resource allocation and improve the recovery of clinical service capacity.

\vspace{1cm}

\hspace{-1.3cm}Keywords: Clinical Engineering, ETL, Multicriteria optimization, Resource allocation, Medical equipment management.